\begin{ouverture}[courbe]{3}{Calcul numérique et littéral}{Partie 1 · Nombres et calculs}
{Puissances · Racines · Développement · Factorisation}
\objectif{Calculer exactement avec puissances et racines.}
\objectif{Choisir une forme adaptée à un objectif.}
\objectif{Manipuler une expression fractionnaire avec ses conditions.}
\prerequis{Priorités opératoires et distributivité.}
\end{ouverture}
\pagestyle{otMain}\essentiel\fontsize{11}{15}\selectfont
\essSection{1}{Puissances et racines}
\begin{propriete}Pour $a\ne0$ et $m,n\in\mathbb Z$ : $a^ma^n=a^{m+n}$, $a^m/a^n=a^{m-n}$, $(a^m)^n=a^{mn}$ et $a^{-n}=1/a^n$. L’écriture scientifique d’un nombre positif est $c\times10^p$ avec $1\le c<10$.\end{propriete}
\begin{definition}Pour $a\ge0$, $\sqrt a$ est l’unique réel positif dont le carré vaut $a$. Ainsi $(\sqrt a)^2=a$ et, pour tout réel $u$, $\sqrt{u^2}=|u|$.\end{definition}
\begin{propriete}Pour $a,b\ge0$, $\sqrt{ab}=\sqrt a\sqrt b$. Pour $a\ge0$ et $b>0$, $\sqrt{a/b}=\sqrt a/\sqrt b$.\end{propriete}
\begin{exemple}$\sqrt{72}=\sqrt{36\times2}=6\sqrt2$ ; $\sqrt{(-5)^2}=5$, et non $-5$.\end{exemple}
\begin{savoirfaire}[Éviter les erreurs de signe]Les parenthèses comptent : $(-3)^2=9$, mais $-3^2=-9$. Pour une écriture fractionnaire, simplifier des facteurs, jamais des termes isolés d’une somme.\end{savoirfaire}
\clearpage
\essSection{2}{Transformer une expression}
\begin{propriete}
\[(a+b)^2=a^2+2ab+b^2,\qquad(a-b)^2=a^2-2ab+b^2,\]
\[(a-b)(a+b)=a^2-b^2.\]
Développer transforme un produit en somme ; factoriser fait l’inverse.
\end{propriete}
\begin{savoirfaire}[Choisir la forme utile]
Pour calculer une valeur ou lire un coefficient, une forme développée peut être pratique. Pour résoudre un produit nul ou étudier un signe, préférer une forme factorisée. Par exemple $x^2-6x+8=(x-2)(x-4)$.
\end{savoirfaire}
\begin{exemple}[Une expression fractionnaire]
$\dfrac{x^2-9}{x-3}=\dfrac{(x-3)(x+3)}{x-3}=x+3$, seulement pour $x\ne3$. La forme simplifiée ne conserve pas à elle seule l’information sur le domaine initial.
\end{exemple}
\begin{savoirfaire}[Isoler une variable]
Si $V=\pi r^2h$ avec $r>0$, alors $h=V/(\pi r^2)$. Écrire les conditions avant toute division. Les unités peuvent servir à contrôler la formule.
\end{savoirfaire}
\begin{attention}$(a+b)^2$ ne vaut pas $a^2+b^2$ en général. Le terme $2ab$ représente aussi, dans un carré découpé, les deux rectangles de côtés $a$ et $b$.\end{attention}
\clearpage\rubrique{Découverte}
\begin{blocLivre}{ex-3-1}
\begin{exercice}[1]{Puissances et parenthèses}
Calculer $(-2)^4$, $-2^4$, $3^{-2}$ et $5^0$.
\end{exercice}
\end{blocLivre}
\begin{blocLivre}{ex-3-2}
\begin{exercice}[1]{Racines simplifiées}
Simplifier $\sqrt{50}$, $\sqrt{108}$ et $\sqrt{16/25}$.
\end{exercice}
\end{blocLivre}
\begin{blocLivre}{ex-3-3}
\begin{exercice}[1]{Écriture scientifique}
Calculer $A=\dfrac{(3\times10^5)(4\times10^{-2})}{6\times10^2}$ et donner une écriture scientifique.
\end{exercice}
\end{blocLivre}
\begin{blocLivre}{ex-3-4}
\begin{exercice}[1]{Développer et réduire}
Développer $A=3(x-2)-2(2x+5)$ et $B=(x+4)(2x-1)$.
\end{exercice}
\end{blocLivre}
\begin{blocLivre}{ex-3-5}
\begin{exercice}[1]{Reconnaître une identité}
Développer $(2x-3)^2$, puis factoriser $x^2-25$ et $x^2+6x+9$.
\end{exercice}
\end{blocLivre}

\clearpage\rubrique{Entraînement}
\begin{blocLivre}{ex-3-6}
\begin{exercice}[2]{Un facteur commun}
Factoriser $A=5x(x-1)+3(x-1)$ et $B=(x+2)^2-4(x+2)$.
\end{exercice}
\end{blocLivre}
\begin{blocLivre}{ex-3-7}
\begin{exercice}[2]{Une simplification conditionnelle}
Simplifier $\dfrac{x^2-4}{x+2}$ en précisant le domaine. Calculer la valeur pour $x=3$.
\end{exercice}
\end{blocLivre}
\begin{blocLivre}{ex-3-8}
\begin{exercice}[2]{Addition de quotients}
Pour $x\ne0$, réduire $\dfrac{2}{x}+\dfrac{x-1}{3x}$ à un seul quotient.
\end{exercice}
\end{blocLivre}
\begin{blocLivre}{ex-3-9}
\begin{exercice}[2]{Isoler une grandeur}
\begin{questions}
\item Isoler $t$ dans $d=vt$ pour $v>0$.
\item Isoler $r$ dans $S=\pi r^2$ pour $r\ge0$ et $S\ge0$.
\item Isoler $y$ dans $3x-2y=7$.
\end{questions}
\end{exercice}
\end{blocLivre}
\begin{blocLivre}{ex-3-10}
\begin{exercice}[2]{Comparer par différence ou par quotient}
\begin{questions}
\item Pour $x\ge0$, comparer $A=(x+1)^2$ et $B=x^2+1$. Quand sont-ils égaux ?
\item Comparer $\dfrac79$ et $\dfrac9{11}$ en calculant leur quotient.
\end{questions}
\end{exercice}
\end{blocLivre}

\clearpage\rubrique{Problèmes}
\begin{blocLivre}{ex-3-11}
\begin{exercice}[3]{Un programme de calcul}
Choisir un nombre, lui ajouter 3, élever au carré, puis retirer neuf. Exprimer le résultat en fonction de $x$ et montrer qu’il est divisible par 4 lorsque $x$ est un entier pair.
\end{exercice}
\end{blocLivre}
\begin{blocLivre}{ex-3-12}
\begin{exercice}[3]{Calculer astucieusement}
Calculer exactement $1001^2-999^2$ sans développer les deux carrés. Proposer une règle générale pour $(n+1)^2-(n-1)^2$.
\end{exercice}
\end{blocLivre}
\begin{blocLivre}{ex-3-13}
\begin{exercice}[3]{Un carré encadré}
Un carré de côté $x>0$ est entouré d’une bordure de largeur 2. Exprimer l’aire de la bordure et la factoriser. La bordure a une aire de 56 : trouver $x$.
\label{fig-ex-3-13}
\figureGeo{bordure}
\end{exercice}
\end{blocLivre}
\begin{blocLivre}{ex-3-14}
\begin{exercice}[3]{Une identité trompeuse}
Un élève écrit $\dfrac{x+4}{x+2}=\dfrac42=2$. Trouver toutes les valeurs de $x$ pour lesquelles son résultat est malgré tout correct.
\end{exercice}
\end{blocLivre}
\begin{blocLivre}{ex-3-15}
\begin{exercice}[3]{Deux modèles équivalents}
Une aire est décrite par $A(x)=(x+5)^2-(x-1)^2$. Trouver une forme simple, puis déterminer $x\ge1$ pour lequel $A(x)=72$. Expliquer pourquoi l’équation est du premier degré.
\end{exercice}
\end{blocLivre}

\clearpage
